\documentclass{article}
\usepackage{iclr2027_conference,times}
\usepackage[T1]{fontenc}
\usepackage{amsmath,amssymb,amsthm,mathtools,booktabs,graphicx,microtype,xcolor}
\usepackage[hidelinks]{hyperref}
\usepackage{etoolbox}
\hypersetup{pdftitle={The Causal Recursion Dividend},pdfauthor={Anonymous authors}}
\newif\ifresearchdraft
\researchdrafttrue
\ifresearchdraft
\makeatletter
\patchcmd{\@maketitle}{Paper under double-blind review}{Research draft; not yet submitted}{}{}
\makeatother
\fi
\newtheorem{theorem}{Theorem}
\newtheorem{proposition}[theorem]{Proposition}
\newtheorem{corollary}[theorem]{Corollary}
\newtheorem{lemma}[theorem]{Lemma}
\theoremstyle{definition}
\newtheorem{definition}{Definition}
\newcommand{\E}{\mathbb E}
\newcommand{\Var}{\operatorname{Var}}
\newcommand{\TV}{\operatorname{TV}}
\newcommand{\KL}{\operatorname{KL}}
\newcommand{\kl}{\operatorname{kl}}
\newcommand{\R}{\mathbb R}
\newcommand{\ind}{\mathbf 1}
\newcommand{\eps}{\varepsilon}
\newcommand{\cU}{\mathcal U}
\newcommand{\cH}{\mathcal H}
\newcommand{\cS}{\mathcal S}
\newcommand{\cP}{\mathcal P}
\newcommand{\authorR}{\mathsf R}
\newcommand{\authorF}{\mathsf F}
\newcommand{\NominalContrast}{$0.2412$}
\newcommand{\RareInterval}{$[0.1971,0.2667]$}
\newcommand{\CoarseInterval}{$[-0.8979,1.0000]$}
\newcommand{\NaiveCoverage}{19.28}

\title{The Causal Recursion Dividend:\\Identifying Improvement of the Improver}
\author{Anonymous authors}
\begin{document}
\maketitle
\ifresearchdraft\lhead{Research draft in ICLR 2027 format; not submitted}\fi
\begin{abstract}
Increasing task performance does not establish that an improvement mechanism became better, or that recursively inheriting it was beneficial. We separate transferable mechanism gain from a comparator-specific \emph{causal recursion dividend}: downstream improvement capability after recursive development minus that after resource-matched fixed-author development. We formulate controlled fork-and-transplant audits and characterize their identification assumptions. For imperfect transplants, finite-state robust evaluation distinguishes independent transition uncertainty from shared history and persistent adapters. In an endpoint-only audit with contrast distortion bounded by $b$, uniformly detecting a true dividend $\gamma$ is impossible when $\gamma\leq2b$; above the boundary, the local minimax sample size scales as $(\gamma-2b)^{-2}\log(1/\delta)$. Development-level confidence bounds avoid treating dependent descendants as independent replications. Exact systems exhibit positive, zero, and negative inheritance effects. Two executable Boolean-repair studies then separate mechanism gains from inheritance gains. A rule-table confirmation gives practical equivalence within one percentage point. An independently confirmed conditional-tree study, with 4,096 development seeds per task family, gives negative dividends of $-1.50$, $-1.34$, and $-1.66$ percentage points, with simultaneous finite-sample intervals excluding zero, although both development methods improve over the seed. These are bounded-system findings, not claims about general-purpose AI. The framework makes claims about improvement of the improver explicit, testable, and sensitive to the declared causal interface.
\end{abstract}

\section{Introduction}
A self-modifying system can improve because it acquires a stronger task solver, a better proposal mechanism, a useful archive, or simply more opportunities to search. These explanations are not interchangeable. Even evidence that the final proposal mechanism outperforms its seed does not establish that \emph{letting improved mechanisms author later modifications} was responsible: a fixed external proposer may discover equally useful mechanisms.

This distinction is increasingly relevant to systems that modify improvement scaffolding. STOP evolves an improver through self-application \citep{stop}; Darwin G\"odel Machine maintains an evolving archive of coding agents \citep{dgm}; Huxley--G\"odel Machine explicitly considers the productivity of descendants \citep{hgm}; and Hyperagents evaluates improvement potential and cross-task transfer \citep{hyperagents}. These works already motivate evaluating improvers and using fixed-mechanism controls. Our goal is not to rename those evaluations. It is to state precisely which causal claim an experiment identifies, which component separations that claim assumes, and how much ambiguity remains when those separations are imperfect.

We distinguish two effects. A \emph{transferable mechanism gain} compares frozen improvers from standardized starting conditions. The \emph{recursion dividend} compares development regimes that differ in whether selected improvers take over subsequent authorship. The nonrecursive comparator is allowed to discover and retain better improvers. It is not restricted to leaving the seed unchanged.

Our central result separates experimental noise from intervention ambiguity. In a bounded terminal-audit model, the gap between a true effect and twice the allowed contrast distortion determines whether uniform certification is possible at all. A finite-state robust audit refines crude worst-case distortion bounds using the actual locations and continuation values of possible mismatches. Both conclusions require an explicit causal interface; more samples cannot replace it.

The mathematical ingredients include robust dynamic programming, bounded-mean concentration, and two-point testing. We identify these established tools rather than presenting them as new general methods. The contribution is their comparator-specific formulation and joint consequences for recursively inherited improvement mechanisms, together with an executable falsification-oriented audit. All proofs and counterexamples are provided. Our CPU experiments demonstrate controlled inheritance effects, practical equivalence in a rule-table setting, and a confirmed negative dividend despite positive mechanism gains in a conditional-program setting. These findings validate attribution rather than a new superior optimization algorithm.

\paragraph{Relation to existing theory.}
Learning mutation policies is already studied in evolutionary optimization \citep{mossina}; neither search-policy adaptation nor symbolic policy representation is our novelty. Robust-testing overlap and least-favorable distributions are established foundations \citep{chen,vallinayagam}. Our endpoint theorem specializes that geometry to intervention distortion. The contribution is the declared authorship intervention, its inferential and interface limitations, and falsification-oriented executable audits, not new general concentration or robust-control machinery.

\section{What effect is being measured?}\label{sec:setup}
A system state contains a task solver $A$, an improvement mechanism $M$, and external history $H$. The experiment declares which code, parameters, and working state belong to each component. The evaluator and budget counter are protected. For task $x$, let $r_x(A)\in[0,1]$ denote externally specified utility. Let $\Psi_B(A,M,H;\xi)$ be the \emph{released} agent after the frozen mechanism $M$ operates with budget $B$ and randomness $\xi$. The mechanism may update ordinary working memory under fixed rules but does not rewrite its own implementation during this audit.

For a prespecified distribution $\mu$ over $(x,A,H)$, define
\begin{align}
 V_B(M;\mu)&=\E_{(x,A,H)\sim\mu,\xi}[r_x(\Psi_B(A,M,H;\xi))],\label{eq:V}\\
 D_B(M;\mu)&=V_B(M;\mu)-V_B(M_0;\mu).\label{eq:D}
\end{align}
The released agent is selected using only information available to the procedure. An oracle maximum over the latent true quality of all descendants is a different estimand. Our program experiment uses exact within-task evaluation; this simplifies selection rather than silently granting one method an oracle.

\paragraph{The development intervention.}
Both regimes begin from the same distribution of full initial states and have the same editing primitives, candidate language, evaluator, selection rule, and development allowance $C$. In \emph{recursive development} $\authorR$, a selected mechanism becomes the proposer of subsequent edits. In \emph{fixed-author development} $\authorF$, the original proposer $M_0$ remains active, but the evolving incumbent mechanism remains editable and can be retained as the final output. Subsequent histories may differ: they are effects of the intervention, not post-treatment quantities to equalize. For selected outputs $M_C^R,M_C^F$, define
\begin{equation}
 \Gamma_{C,B}(\mu)=\E_{\mathrm{development}}[V_B(M_C^R;\mu)-V_B(M_C^F;\mu)].\label{eq:Gamma}
\end{equation}
Thus a positive $D_B$ need not imply a positive $\Gamma_{C,B}$. A positive $\Gamma_{C,B}$ need not imply either developed mechanism beats its seed. Every statement is relative to $\mu,C,B$, the declared component boundary, and the fixed-author comparator.

\paragraph{A semantic, not syntactic, distinction.}
A fixed interpreter can simulate a computable self-modifying process by storing the changing program as data. Consequently, no advantage over \emph{all fixed programs} follows from calling a program recursive. Our comparison freezes proposer \emph{behavior}, not merely the interpreter's source text. Primitive-call counts are matched in the experiments; equal wall time or instruction counts are not assumed.

\section{Identification and its limits}\label{sec:identification}
\begin{proposition}[Native-summary nonidentification]\label{prop:nonid}
Observations consisting only of native performance-and-lineage pairs can fail to identify even the sign of $D_B$, including in deterministic finite systems.
\end{proposition}
To see this, the following two worlds agree on native observations $(A_0,M_0)$ and $(A_1,M_1)$ but reverse both controlled mechanism effects:
\begin{equation}
\begin{array}{c|cc|cc}
 &\multicolumn{2}{c|}{\text{World I}}&\multicolumn{2}{c}{\text{World II}}\\
 &M_0&M_1&M_0&M_1\\\hline
 A_0&.4&.7&.4&.3\\
 A_1&.5&.8&.9&.8
\end{array}\label{eq:worlds}
\end{equation}
Native performance rises from $.4$ to $.8$ in either world; a component swap yields $+.3$ or $-.1$. This elementary counterexample concerns the stated observation model, not an investigator with a complete executable environment model and unrestricted counterfactual queries.

\paragraph{Identification by forks.}
Generate independent development pairs, randomizing exogenous seeds and cloning their initial states. Apply $\authorR$ and $\authorF$, freeze their selected mechanisms, then evaluate both on common fresh draws from $\mu$. With faithful interfaces, branch consistency, and no cross-branch interference, the average paired released-score difference is unbiased for \eqref{eq:Gamma}. Paired randomness is a variance-reduction device, not an identification assumption. Checkpoint selection may use development data; it must not silently reuse the final audit as development feedback.

\begin{proposition}[Where inheritance contributes]\label{prop:decomp}
Let $Z_t$ be a sufficient finite development state, $P_t^R,P_t^F$ the two regime kernels, $g_B(Z_C)=V_B(M(Z_C);\mu)$, $d_t^R$ the recursive occupancy, and $V_t^F$ the fixed-author continuation value. For a common initial law,
\begin{equation}
 \Gamma_{C,B}=\sum_{t=0}^{C-1} d_t^R(P_t^R-P_t^F)V_{t+1}^F.\label{eq:decomposition}
\end{equation}
\end{proposition}
This is a finite-horizon performance-difference identity, not a new general lemma \citep{trpo}. It isolates a substantive distinction: an inherited mechanism's immediate task score is not its ability to propose states with higher \emph{future fixed-author continuation value}. Even monotone selection of better task improvers can yield a negative dividend. Section~\ref{sec:experiments} constructs this case explicitly.

\section{Imperfect transplants: which effects are identifiable?}\label{sec:robust}
Swapping a mechanism may also alter an archive adapter, a representation, or an undeclared shared state. We do not infer the size of this mismatch from a nominal performance gain. Instead, the audit specifies a justified uncertainty set for the intended intervention. Uncertainty radii in our controlled examples are structural inputs, not learned calibration guarantees for an arbitrary agent.

For a fixed development pair, branch $a\in\{R,F\}$ has nominal audit kernels $\widetilde K_t^a$ on a sufficient finite state space $\cS$, known initial law $\nu_a$, and terminal reward $r_a\in[0,1]$. Each intended transition row belongs to
\begin{equation}
 \cU_{t,s}^a=\{q\in\Delta(\cS):\TV(q,\widetilde K_t^a(s,\cdot))\leq\eps_{t,s}^a,
 \ \operatorname{supp}(q)\subseteq E_{t,s}^a\}.\label{eq:uncertainty}
\end{equation}
The nominal row satisfies the declared support. $E$ encodes sound structural zeros, not unseen transitions in a small sample. First assume uncertainty is independent across time, state, and branch: a \emph{rectangular} model.

\begin{theorem}[Sharp rectangular audit interval]\label{thm:rectangular}
Set $L_h^a=U_h^a=r_a$, and recursively define
\begin{equation}
 L_t^a(s)=\min_{q\in\cU_{t,s}^a}q^\top L_{t+1}^a,
 \qquad U_t^a(s)=\max_{q\in\cU_{t,s}^a}q^\top U_{t+1}^a.\label{eq:bellman}
\end{equation}
The sharp identified interval for the intended paired effect is
\begin{equation}
 \left[\nu_R^\top L_0^R-\nu_F^\top U_0^F,
       \nu_R^\top U_0^R-\nu_F^\top L_0^F\right].\label{eq:sharp}
\end{equation}
Both endpoints are attained by admissible kernels, and every intermediate value is attainable. Each TV-row optimization is solved by transferring probability mass from larger to smaller continuation values (reversed for maximization).
\end{theorem}
Equation~\eqref{eq:bellman} uses established robust dynamic programming \citep{nilim}. Its audit consequence is that mismatch matters through \emph{reachable, decision-relevant continuation differences}, not only its maximum magnitude. A crude coupling bound uses $\beta_a=1-\prod_t(1-\eps_t^a)$ when every row at time $t$ has radius at most $\eps_t^a$, and expands the nominal contrast by $\beta_R+\beta_F$. It can be much wider than \eqref{eq:sharp}.

\paragraph{Shared uncertainty is different.}
Independent branch balls are generally only an outer bound when an unknown adapter or history parameter is shared. Suppose the \emph{same} initial history law $q$ is within TV distance $\eps$ of $q_0$, and conditional branch values $v_R(h),v_F(h)$ are known. The sharp interval is
\begin{equation}
 \left[\min_{\TV(q,q_0)\leq\eps}q^\top(v_R-v_F),
       \max_{\TV(q,q_0)\leq\eps}q^\top(v_R-v_F)\right].\label{eq:shared}
\end{equation}
If $v_R-v_F$ is constant, the causal contrast is exactly identified even when either branch's value is highly uncertain. In general,
$|q^\top(v_R-v_F)-q_0^\top(v_R-v_F)|\leq\eps\,\operatorname{span}(v_R-v_F)$.
Thus branchwise uncertainty need not equal attribution uncertainty. These shared-law and rectangular cases cover different causal assumptions; one must not report rectangular sharpness for an entangled unknown parameter.

\paragraph{Persistent adapters.}
A single unknown adapter reused across an audit is not an independently uncertain transition at every step. Proposition~\ref{prop:persistent} gives the exact finite-adapter identified set. A two-step example has contrast $+.2$ for every persistent adapter, but its rowwise relaxation has interval $[-.2,.2]$. The negative endpoint is caused by switching adapters mid-audit and is inadmissible in the persistent model. Thus a conservative relaxation can erase sign identification without demonstrating a genuinely adverse interface.

\section{The irreducible detection boundary}\label{sec:boundary}
Sharp finite-state intervals are useful when transition structure is known. A simpler endpoint model isolates what sampling alone can achieve. Let ideal released rewards be Bernoulli with means $\theta_R,\theta_F$, true dividend $\theta=\theta_R-\theta_F$, and observed means $p_R,p_F$. Assume
\begin{equation}
 |p_R-\theta_R|\leq b/2,\quad |p_F-\theta_F|\leq b/2.\label{eq:endpoint}
\end{equation}
Here $b$ bounds \emph{contrast} distortion, not one-branch distortion. An endpoint-only audit obtains $n$ IID observed pairs and cannot acquire additional structural information.

\begin{theorem}[Minimax separation under transplant distortion]\label{thm:gap}
Let $0\leq b\leq1$, $0<\gamma\leq1$. Consider testing $H_0:\theta\leq0$ against $H_1:\theta\geq\gamma$, with both errors at most $0<\delta<1/2$ uniformly over \eqref{eq:endpoint}. If $0<\gamma\leq2b$, no such test exists for any sample size. If $0\leq b<1/8$ and $2b<\gamma\leq1/4$, write $g=\gamma-2b$. Thresholding the observed mean contrast at $\gamma/2$ suffices when
\begin{equation}
 n\geq\frac{8\log(1/\delta)}{g^2}.
\end{equation}
Conversely, even with independent arms, every valid test requires
\begin{equation}
 n\geq\frac{\kl(1-\delta,\delta)}{4g^2},\qquad
 \kl(1-\delta,\delta)=(1-2\delta)\log\frac{1-\delta}{\delta}.\label{eq:lower}
\end{equation}
\end{theorem}
The lower bound constructs overlapping observation laws below the boundary and uses Bernoulli KL above it. This is a specialization of robust-testing reasoning, not a new general contamination theory \citep{chen,vallinayagam}. The factor two has a precise meaning: an \emph{observed} positive contrast exceeding $b$ can certify a positive intended effect, but uniform power for all intended effects at least $\gamma$ requires $\gamma>2b$. Sampling error vanishes with $n$; the intervention ambiguity does not. The result concerns the endpoint observation model, not every possible structured intervention experiment.

\section{Independent development, not independent descendants}\label{sec:stats}
Let development pair $\mathcal D_i$ produce a frozen pair of mechanisms. Conditional on $\mathcal D_i$, let $W_{ij}\in[-1,1]$ be fresh IID paired task-audit differences with mean $\Delta_i$. Set $Z_i=m^{-1}\sum_{j=1}^mW_{ij}$. Independent development pairs make the $Z_i$ IID; individual descendants within a pair are not marginally independent.

\begin{proposition}[Finite-sample fork-and-transplant certificate]\label{prop:certificate}
For $J$ prespecified contrasts and $n$ independent development pairs, let $\bar Z$ and $s_Z^2$ be the sample mean and unbiased sample variance. With probability at least $1-\delta$, simultaneously for all $J$ effects, a valid interval is $[\bar Z-r-b,\bar Z+r+b]\cap[-1,1]$, where a uniform bound $b$ controls contrast distortion and either
\begin{align}
 r_H&=\sqrt{2\log(2J/\delta)/n},\label{eq:hoeffding}\\
 r_{EB}&=\sqrt{2s_Z^2\log(4J/\delta)/n}
          +\frac{14\log(4J/\delta)}{3(n-1)},\quad n\geq2.\label{eq:eb}
\end{align}
The choice of bound is specified in advance. A lower endpoint above a practical threshold $\eta$ certifies an effect larger than $\eta$.
\end{proposition}
The empirical Bernstein bound is obtained by rescaling \citet[Theorem~4]{maurer}; it is not new. The statement is fixed-sample, not optional-stopping-valid. Time-uniform inference needs a confidence-sequence construction \citep{howard} and an audit-data protocol that still controls adaptive model selection. A uniform $b$ must be justified separately; substituting a fitted, unvalidated mismatch estimate does not preserve coverage.

\begin{proposition}[Replication floor and allocation]\label{prop:variance}
If $\tau^2=\Var(\Delta_i)$ and $\sigma^2=\E[\Var(W_{ij}\mid\mathcal D_i)]$, then
\begin{equation}
 \Var(\bar Z)=\frac{\tau^2+\sigma^2/m}{n}.\label{eq:variance}
\end{equation}
For cost $n(c_d+mc_a)$ and positive $\tau^2,c_d,c_a$, the unconstrained continuous optimal inner replication is $m^*=\max\{1,\sqrt{c_d\sigma^2/(c_a\tau^2)}\}$, subject to feasibility and integer rounding. Even infinitely many descendants cannot eliminate the $\tau^2/n$ term; a two-point subfamily requires $\Omega(\eps^{-2})$ independent development runs for constant-confidence accuracy $\eps$.
\end{proposition}
This is the ordinary nested-variance identity with an important audit interpretation. Sampling 10,000 descendants of one discovered improver estimates that improver well; it does not establish what the development algorithm reliably discovers.

\section{Experiments}\label{sec:experiments}
We run three controlled suites and two executable improver representations on Boolean-program repair. All use synthetic tasks, single-process CPU computation, and no model training or API calls. Protocols, random seeds, per-run outcomes, source hashes, and exact-value checks accompany the paper. Initial and independent-confirmation samples are kept separate.

\subsection{Exact systems and attribution sensitivity}
Four mechanism states have single-proposal success probabilities $(.02,.04,.20,.30)$ and downstream value $g_B(j)=1-(1-p_j)^B$. Development proposes a mechanism state and retains it only if better. A fixed author uses the seed proposal distribution at every step; the recursive author uses the current state's distribution. Three declared proposal matrices produce positive, zero, and negative dividends. Both regimes can acquire better improvers in all three cases. Exact matrix products and 4,096 paired simulations per setting agree, and the telescoping identity holds to numerical precision. The first generation has zero dividend because both authors initially coincide.

\begin{figure}[t]
\centering\includegraphics[width=.89\linewidth]{figures/exact_dividend.pdf}
\caption{Exact finite-state recursion dividends at downstream budget $B=4$. Retaining stronger task improvers can help, harm, or leave unchanged the value of recursive authorship. The fixed-author comparator can improve in every system.}
\label{fig:exact}
\end{figure}

For imperfect transplants, an eight-step chain has a 5\% rare-state block and a 95\% clean block. Success hazards are $.06$ and $.02$ in the two branches; nominal contrast is \NominalContrast. Equal maximum row radii yield markedly different intervals depending on whether uncertainty affects only the rare block or all states. At $\eps=.10$, the rare-block sharp interval is \RareInterval, while the uniform coupling interval is \CoarseInterval. We verify each sharp endpoint by evaluating its returned extremal kernels. A separate shared-history example preserves an exact contrast of $.1$ despite substantial uncertainty in either branch.

\begin{figure}[t]
\centering\includegraphics[width=.89\linewidth]{figures/sensitivity.pdf}
\caption{Attribution depends on uncertainty structure. Solid boundaries are sharp with uncertainty confined to a rare block; dashed boundaries allow it globally. The coarse uniform coupling bound (dotted) discards reachability. No sampling uncertainty is included in this exact calculation.}
\label{fig:sensitivity}
\end{figure}

\subsection{Descendant dependence and endpoint testing}
A cluster model uses 16 independent development runs, conditional mean $\Delta_i\in\{-.4,+.4\}$, and $m$ descendants per run. Across 10,000 repetitions, incorrectly treating all descendants as independent drives nominal 95\% Hoeffding coverage down to \NaiveCoverage\% at $m=1,024$; the development-level bound remains conservative. Empirical variance follows \eqref{eq:variance}. We also simulate the null/alternative observation laws from Theorem~\ref{thm:gap}: below the separation boundary the same observed law persists at every sample size, whereas errors shrink above it. These are calibration and falsification tests, not demonstrations of general AI capabilities.

\subsection{Executable Boolean-program repair}
Task agents are feed-forward circuits with six inputs and sixteen gates, using AND, OR, XOR, and NAND. A target circuit is sampled from one of three operation mixtures (mixed, XOR-heavy, or logic-heavy); a parent is formed by corrupting five active genes. Trivial targets and functionally correct parents are rejected. All 64 inputs fit in one machine word, making exact functional evaluation inexpensive.

The improver is an eight-integer rule table controlling edit-position weights, edit types, active-cone preference, and acceptance of neutral task edits. A fixed typed mutation interpreter applies this rule table both to circuit genomes and to the rule table itself. Recursive development uses the selected rule table as the next author; fixed-author development always uses the seed rule table to edit the evolving incumbent. A random-search baseline proposes entire rule tables uniformly. This is bounded self-referential search, not an unrestricted code-writing agent.

Each development run evaluates its seed and $C$ proposals on six training tasks with repair budget 24, retaining strict training-score improvements. Training tasks and search randomness are fixed within a run and independent across runs. At audit, policies are frozen and evaluated on 32 fresh task--parent pairs. Search outcomes are paired across methods by common random numbers. The full 64-input evaluator is available during within-task selection, so generalization is across new tasks and parents, not unseen truth-table inputs. No audit scores feed back into policy selection.

The initial study uses 512 development runs per family and reports every cell of $C\in\{16,64,128\}$ and $B\in\{8,32,128\}$. After observing substantial mechanism gains but small dividends, we freeze a separate confirmation protocol: 4,096 new runs per family, $C=128,B=32$, a one-percentage-point practical equivalence band, and empirical Bernstein intervals simultaneous across the three primary dividends. The equivalence margin and variance-sensitive method were chosen after the initial study but before any confirmation outcomes; this is independent confirmation, not retrospective preregistration. Nine seed-relative gains form a separate 95\% simultaneous confidence family.

\begin{table}[t]
\centering
\caption{Independent confirmation: effects in percentage points of correct truth-table outputs. Intervals for $\Gamma$ use the prespecified finite-sample empirical Bernstein bound, simultaneous across the three families. $D_R,D_F$ are point estimates; their separate simultaneous intervals appear in Appendix~\ref{app:results}.}
\label{tab:confirmation}
\begin{tabular}{lrrrr}
\toprule
Family & $D_R$ & $D_F$ & $\Gamma$ & 95\% interval for $\Gamma$ \\
\midrule
Mixed & 5.55 & 5.61 & -0.065 & $[-0.841,+0.711]$ \\
Xor & 4.08 & 4.10 & -0.021 & $[-0.771,+0.730]$ \\
Logic & 6.21 & 6.15 & +0.058 & $[-0.725,+0.842]$ \\
\bottomrule
\end{tabular}

\end{table}

Table~\ref{tab:confirmation} separates two claims that an ordinary progress curve conflates. Both developed mechanisms outperform their seed, but all three primary dividend intervals lie inside $[-1,1]$ percentage points. This establishes practical equivalence at the declared tolerance in this benchmark, not exact equality and not a universal absence of recursion benefits. Uniform random policy search is a stronger comparator here than either local development rule; the complete results retain this unfavorable baseline. The benchmark therefore validates attribution, not a superior RSI algorithm.

Each regime uses $129$ policy evaluations and $129\cdot6\cdot25=19,350$ development circuit evaluations per run. A final $B=32$ audit uses $32\cdot33=1,056$ circuit evaluations per selected policy. We charge all candidate evaluations, including rejected and no-op proposals. The protected interpreter's control-flow costs differ; we match evaluator counts and report measured CPU wall time separately. A claim of identical instruction-level compute would be stronger than this experiment supports.

\subsection{Conditional programs: positive mechanism gains, negative inheritance}\label{sec:trees}
A second improver is an executable depth-two decision tree, not a fixed weight table. Three internal nodes branch on elapsed budget, stagnation, score, active-gene fraction, or normalized genotype value. Four leaves select a typed coordinate restriction and reset/increment/decrement operation. The same protected editing primitive acts on circuit genomes and the tree's ten-integer syntax. Recursive development installs the selected tree as the next author; fixed-author development retains the uniform seed author while editing and retaining the same kind of incumbent. Both use strict training improvement. A whole-tree uniform random-search baseline is also evaluated.

The exploratory study uses 2,048 development seeds per family and checkpoints $C\in\{16,64\}$, $B\in\{16,64\}$. Its mixed-family negative dividend motivates a separate confirmation: 4,096 new seeds per family, $C=64,B=64$, unchanged grammar, seed, task generator, and selection rule. This protocol was fixed before confirmation outcomes. Each policy is developed on six tasks at repair budget 24 and audited on 24 fresh tasks. The three primary dividend intervals use the same finite-sample empirical Bernstein construction with $J=3$; nine seed gains form a separate $J=9$ confidence family. No confirmation seeds are pooled with exploration.

\begin{table}[t]
\centering
\caption{Conditional-tree independent confirmation: percentage-point effects on fresh tasks. Both local development methods improve over the seed, but all three simultaneous dividend intervals are strictly negative.}
\label{tab:trees}
\begin{tabular}{lrrr}
\toprule
Family & $D_R$ & $D_F$ & $\Gamma$ [simultaneous 95\% interval] \\
\midrule
mixed & 3.28 & 4.78 & -1.50 [-2.35, -0.65] \\
xor & 2.89 & 4.23 & -1.34 [-2.17, -0.51] \\
logic & 3.75 & 5.41 & -1.66 [-2.53, -0.80] \\
\bottomrule
\end{tabular}

\end{table}

Table~\ref{tab:trees} establishes a negative inheritance effect under this declared comparator, rather than merely failing to detect a positive effect. Every seed-relative gain has a positive lower bound in its separate confidence family (Appendix~\ref{app:trees}). Uniform whole-tree search is stronger in mean than both local methods and is retained in the results. Post-hoc traces show more condition edits, fewer leaf edits, and fewer accepted improvements under inherited authorship. These are diagnostics, not an identified mediation mechanism. An additional all-permutation meta-type intervention leaves the fixed-author comparator pointwise unchanged (Proposition~\ref{prop:adapterinvariance}); it detects no adapter benefit, with intervals too wide to establish invariance (Appendix~\ref{app:adapterdiagnostic}).

Each method uses $65\cdot6\cdot25=9,750$ development circuit evaluations per run; each frozen-policy audit uses $24\cdot65=1,560$. Because the tree observes elapsed budget fraction, each downstream budget is restarted separately, not read from a longer-budget prefix. The $C=16$ exploratory checkpoint belongs to a $C=64$ schedule. The two representation studies differ in budgets and policy semantics, so their difference is not itself an isolated causal effect of representation. Their common lesson is that improving an improver and benefiting from its recursive inheritance are different empirical claims.

\section{Scope, limitations, and conclusion}\label{sec:limitations}
The proposed dividend is relative to a declared comparator and module boundary. Fixed programs can emulate self-modification, and information embedded in code cannot in general be separated uniquely from an algorithm. Our finite-state intervals require known structural uncertainty sets; estimating or validating those sets for realistic agents is a separate problem. Shared hidden parameters can invalidate rectangular sharpness, and undeclared persistent state can invalidate a Markov model. The testing boundary is sharp for the endpoint model stated, not for arbitrary richer observations.

Our experiments use two small interpretable representations within one Boolean-repair domain. They do not validate broad transfer, language-model behavior, or unrestricted recursive improvement. The separation is stronger than a null finding alone: learning a better improver can coexist with a statistically certified disadvantage from inheriting authorship in the declared system. The paper's tools are intended to make such distinctions auditable. Progress, mechanism improvement, and recursive advantage should be reported as different claims, with explicit intervention assumptions and uncertainty appropriate to the development process.

\clearpage
\section*{AI-use disclosure}
Generative AI assisted the formulation, derivation and checking of proofs, implementation, synthetic-experiment design and execution, analysis, and manuscript drafting. Automated tests and executed results are documented in the supplement; they are not a substitute for independent mathematical review. This working draft has not undergone a completed human-author verification pass. Human authors must verify the claims, references, code, and disclosure before submission. No generative model was trained or queried by the experimental algorithms.

\section*{Reproducibility and ethics}
The supplement contains the full source, frozen protocols, aggregate and per-development-run results, tests, and all proofs. The main experiments use no personal data, human participants, external execution privileges, or paid APIs. Neither editable representation can modify the interpreter, evaluator, filesystem, or budget counter. Identifying public-repository links are deliberately absent from this anonymous-format manuscript. A public working repository is not an anonymous supplementary link.

\begin{thebibliography}{12}
\bibitem[Chen et~al.(2016)Chen, Gao, and Ren]{chen}
Mengjie Chen, Chao Gao, and Zhao Ren.
A general decision theory for Huber's $\epsilon$-contamination model.
\emph{Electronic Journal of Statistics}, 10(2):3752--3774, 2016.
\url{https://arxiv.org/abs/1511.04144}.

\bibitem[Howard et~al.(2021)Howard, Ramdas, McAuliffe, and Sekhon]{howard}
Steven R. Howard, Aaditya Ramdas, Jon McAuliffe, and Jasjeet Sekhon.
Time-uniform, nonparametric, nonasymptotic confidence sequences.
\emph{Annals of Statistics}, 49(2):1055--1080, 2021.
\url{https://arxiv.org/abs/1810.08240}.

\bibitem[Maurer and Pontil(2009)]{maurer}
Andreas Maurer and Massimiliano Pontil.
Empirical Bernstein bounds and sample variance penalization.
In \emph{Conference on Learning Theory}, 2009.
\url{https://arxiv.org/abs/0907.3740}.

\bibitem[Nilim and El~Ghaoui(2005)]{nilim}
Arnab Nilim and Laurent El~Ghaoui.
Robust control of Markov decision processes with uncertain transition matrices.
\emph{Operations Research}, 53(5):780--798, 2005.
\url{https://people.eecs.berkeley.edu/~elghaoui/pubs_rob_mdp.html}.

\bibitem[Schulman et~al.(2015)Schulman, Levine, Moritz, Jordan, and Abbeel]{trpo}
John Schulman, Sergey Levine, Philipp Moritz, Michael I. Jordan, and Pieter Abbeel.
Trust region policy optimization.
In \emph{International Conference on Machine Learning}, pp. 1889--1897, 2015.
\url{https://arxiv.org/abs/1502.05477}.

\bibitem[Wang et~al.(2025)Wang, Pi\k{e}kos, Li, Laakom, Chen, Ostaszewski, Zhuge, and Schmidhuber]{hgm}
Wenyi Wang, Piotr Pi\k{e}kos, Li Nanbo, Firas Laakom, Yimeng Chen, Mateusz Ostaszewski, Mingchen Zhuge, and J\"urgen Schmidhuber.
Huxley--G\"odel Machine: Human-level coding agent development by an approximation of the optimal self-improving machine.
\emph{arXiv preprint arXiv:2510.21614}, 2025.
\url{https://arxiv.org/abs/2510.21614}.

\bibitem[Zelikman et~al.(2024)Zelikman, Lorch, Mackey, and Kalai]{stop}
Eric Zelikman, Eliana Lorch, Lester Mackey, and Adam Tauman Kalai.
Self-taught optimizer (STOP): Recursively self-improving code generation.
In \emph{Conference on Language Modeling}, 2024.
\url{https://arxiv.org/abs/2310.02304}.

\bibitem[Zhang et~al.(2025)Zhang, Hu, Lu, Lange, and Clune]{dgm}
Jenny Zhang, Shengran Hu, Cong Lu, Robert Lange, and Jeff Clune.
Darwin G\"odel Machine: Open-ended evolution of self-improving agents.
\emph{arXiv preprint arXiv:2505.22954}, 2025. Version 3, March 2026.
\url{https://arxiv.org/abs/2505.22954}.

\bibitem[Zhang et~al.(2026)Zhang, Zhao, Yang, Foerster, Clune, Jiang, Devlin, and Shavrina]{hyperagents}
Jenny Zhang, Bingchen Zhao, Wannan Yang, Jakob Foerster, Jeff Clune, Minqi Jiang, Sam Devlin, and Tatiana Shavrina.
Hyperagents.
\emph{arXiv preprint arXiv:2603.19461}, 2026.
\url{https://arxiv.org/abs/2603.19461}.

\bibitem[Mossina et~al.(2019)Mossina, Rachelson, and Delahaye]{mossina}
Luca Mossina, Emmanuel Rachelson, and Daniel Delahaye.
A reinforcement learning perspective on the optimal control of mutation probabilities for the (1+1) evolutionary algorithm: First results on the OneMax problem.
\emph{arXiv preprint arXiv:1905.03726}, 2019.
\url{https://arxiv.org/abs/1905.03726}.

\bibitem[Vallinayagam et~al.(2026)Vallinayagam, Pensia, and Jog]{vallinayagam}
Shankar Vallinayagam, Ankit Pensia, and Varun Jog.
On the sample complexity of robust binary hypothesis testing.
\emph{arXiv preprint arXiv:2605.24741}, 2026.
\url{https://arxiv.org/abs/2605.24741}.
\end{thebibliography}

\clearpage
\appendix
\section{Complete proofs}\label{app:proofs}
\subsection{Native-summary nonidentification and semantic emulation}
\begin{proof}[Proof of Proposition~\ref{prop:nonid}]
Fix any $B$ and take a degenerate task distribution. Define the released-score map on the four possible parent--mechanism pairs by either matrix in \eqref{eq:worlds}. A deterministic finite environment can implement this map by storing the pair and emitting its table entry as terminal reward. Let the development record visit $(A_0,M_0)$ followed by $(A_1,M_1)$, with the same names, parent pointers, budgets, and recorded scores in both worlds. Then the entire specified record is identical. On either standardized parent, World I's controlled mechanism contrast is $.3$, while World II's is $-.1$. Any statistic of the identical record has the same value in the two worlds and therefore cannot identify the sign. To embed the construction in a longer horizon, insert deterministic transitions that preserve the recorded pair and score; do not add the unobserved cross-pair interventions. This proves nonidentification for the declared summaries.
\end{proof}

\begin{proposition}[Fixed-interpreter emulation]\label{prop:emulation}
For any finite-budget computable self-modifying process with a specified random-input stream and primitive-call interface, there exists a fixed interpreter whose state includes the process's editable code and whose primitive-call trace is identical for every random-input stream. This does not imply equal machine-instruction or wall-clock cost.
\end{proposition}
\begin{proof}
Represent the current editable program, working memory, program counter, and remaining budget as data in the interpreter state. At each step the interpreter executes the next transition of the represented program. A self-edit changes the represented program data; an external primitive call is forwarded unchanged; a random-input request consumes the same next random symbol. Induction on the finite execution length gives equality of represented state and forwarded primitive trace at every step. The interpreter itself is fixed, but its interpreted proposer behavior need not remain the seed behavior. Interpreter overhead is not controlled by the argument.
\end{proof}

This proposition is a representation limitation, not an argument that recursive inheritance never helps. It explains why our fixed-author baseline restricts proposer semantics and why the random-search baseline provides an additional substantive comparison. Our emulation software test checks equality of the declared reference semantics, not a general compiler-correctness theorem.

\subsection{Identification by controlled development and transplantation}
Write $Y^a(\mathcal D_i,X_{ij},\xi_{ij})$ for the released score obtained by faithfully installing the selected mechanism from regime $a$ in audit case $X_{ij}\sim\mu$. Assume that the branch has no access to the other branch's changing state; that the intended installation and the executed installation agree; and that the sampling law for development is the one appearing in \eqref{eq:Gamma}. Conditional on development,
\[
 \E[Y^R-Y^F\mid\mathcal D_i]
 =V_B(M_C^R;\mu)-V_B(M_C^F;\mu).
\]
Taking expectations over independent development pairs gives $\E Z_i=\Gamma_{C,B}$. The same argument with a fixed seed comparator identifies $D_B$. Coupling $X_{ij}$ or $\xi_{ij}$ between branches changes covariance, not the marginals and hence not this equality. Different branches need not produce identical post-treatment archives. Installation of a common \emph{initial} archive is part of $\mu$; forcing later archives to coincide would instead change the intervention.

This result estimates the value of a development algorithm averaged over its specified randomness. It does not establish the quality of every possible checkpoint or a population outside $\mu$. If a checkpoint is selected using a held-out audit and then reported with an unadjusted confidence interval on the same audit, the unbiasedness argument for a fixed selected mechanism no longer applies directly. The finite-family correction in Proposition~\ref{prop:certificate} covers selection among a fixed prespecified set, not unlimited adaptive creation of new mechanisms using audit results.

\subsection{Development performance difference}
\begin{proof}[Proof of Proposition~\ref{prop:decomp}]
For the fixed-author chain, $V_C^F=g_B$ and $V_t^F=P_t^FV_{t+1}^F$. Under recursive development,
\[
 \E_R[V_{t+1}^F(Z_{t+1})-V_t^F(Z_t)]
 =d_t^R(P_t^R-P_t^F)V_{t+1}^F.
\]
Sum over $t=0,\ldots,C-1$. The left side telescopes to
$\E_R[g_B(Z_C)]-\E_R[V_0^F(Z_0)]$.
Because the two regimes have the same initial law, the second term equals $\E_F[g_B(Z_C)]$. Their difference is \eqref{eq:Gamma}. The state must contain all history needed to make the kernels Markov; an incumbent's score alone will generally not suffice.
\end{proof}

\subsection{TV optimization and rectangular sharpness}
\begin{lemma}[TV mass-transfer optimization]\label{lem:tv}
For a probability vector $p$ on a finite known support $E$, a finite vector $v$, and $\eps\in[0,1]$, the minimum of $q^\top v$ subject to $q\in\Delta(E)$ and $\TV(p,q)\leq\eps$ is attained by moving mass from highest-value donor coordinates to lowest-value receiver coordinates until the TV budget is exhausted or no strict improvement remains. The maximum follows by replacing $v$ with $-v$.
\end{lemma}
\begin{proof}
Write $(q-p)_+$ and $(p-q)_+$ for received and donated mass. They have equal total $t=\TV(p,q)\leq\eps$. Among donors, removing a unit at a larger value cannot be worse than removing it at a smaller value; exchange the donor assignments whenever they violate this order, subject to each donor's available mass $p_i$. Likewise, assigning a received unit to a smaller value cannot be worse than assigning it to a larger value, subject to coordinate capacity $1-p_i$. A donor and receiver can be combined or canceled if the same coordinate appears in both roles, without increasing total moved mass or changing $q$. Consequently an optimum can be arranged with descending-value donors and ascending-value receivers. Exhaust each available donor or receiver before advancing to the next. Stop when donor value does not exceed receiver value, since additional transfers then cannot improve the objective. This is exactly the greedy algorithm. Its output is feasible, and any feasible transfer plan can be transformed to this ordering without increasing the objective, proving optimality. Ties permit multiple optima.
\end{proof}

\begin{proof}[Proof of Theorem~\ref{thm:rectangular}]
Each row set is a nonempty compact convex subset of the simplex. For terminal time $h$, the claimed lower and upper values equal the fixed reward. Suppose the lower value is correct at time $t+1$ for every state, and that one admissible future collection of rows achieves these lower values simultaneously. Such simultaneous attainment holds at time $h$ trivially. For any admissible row $q$ at $(t,s)$ and any admissible continuation, expected reward is at least $q^\top L_{t+1}^a$, hence at least $L_t^a(s)$. Compactness gives a minimizing row. Because the row sets are independent across states and times, choose a minimizing row at every state and combine them with the simultaneously minimizing future collection. This achieves $L_t^a$ at all states. Backward induction proves the lower result and simultaneous attainment; the maximum is identical with inequalities reversed.

The fixed initial law gives the branch extrema $\nu_a^\top L_0^a$ and $\nu_a^\top U_0^a$. Independence of branch uncertainty permits combining the minimizing recursive kernels with maximizing fixed-author kernels, and vice versa. This proves the endpoints in \eqref{eq:sharp}. The full product of row sets is connected and compact. The contrast is a continuous function of its entries (a finite sum of products). Its image in $\R$ is therefore a compact interval containing its minimum and maximum, so all intermediate effects are attained. Lemma~\ref{lem:tv} supplies the stated row optimization.
\end{proof}

\paragraph{Why rectangularity matters.}
If two transition rows must share an unknown parameter, independently choosing their minimizing values can violate the joint constraint. The construction in the proof is then infeasible. The rectangular relaxation still provides an outer interval if it contains the true joint set, but equality and endpoint attainment for the original problem need not follow. Similarly, if hidden persistent state has been omitted, there may not be a single intended Markov row corresponding to the declared state; merely assigning a small row radius does not repair the model.

\begin{corollary}[Shared-history sharpness and invariance]\label{cor:shared}
Under the common-history-law model in \eqref{eq:shared}, that interval is sharp. If $d(h)=v_R(h)-v_F(h)$ is constant on the allowed histories, the identified set is the singleton containing that constant. For any $q$ in the TV ball,
\[
 |(q-q_0)^\top d|\leq\eps(\max_h d(h)-\min_h d(h)).
\]
\end{corollary}
\begin{proof}
Condition on history and apply the law of total expectation to obtain the intended contrast $q^\top d$. The feasible $q$ set is compact and convex and the objective is linear, so the extrema in \eqref{eq:shared} are attained and all intermediate values follow by convex mixing. For the inequality, subtract the constant $\min_h d(h)$ from $d$; $q-q_0$ sums to zero. The sum of its positive entries is $\TV(q,q_0)$, so the absolute signed expectation is bounded by this mass times the range of $d$. A constant $d$ has range zero. The argument also applies to a known common pre-fork Markov process by propagating $d$ backward; optimizing its uncertain kernels requires preserving their shared structure.
\end{proof}

For example, $q_0=(.5,.5)$, $v_R=(.3,.9)$, $v_F=(.2,.8)$, and radius $1$ allow either branch's value to vary over an interval of width $.6$. Nevertheless the intended contrast is always $.1$. Treating the histories as independently uncertain between branches would incorrectly enlarge the sharp interval for the common-history model to $[-.5,.7]$.

\begin{lemma}[Coarse finite-horizon coupling bound]\label{lem:coupling}
If the nominal and intended kernels have rowwise TV distance at most $\eps_t$ at time $t$, share an initial law, and have the same terminal reward in $[0,1]$, their terminal expected rewards differ by at most $1-\prod_{t=0}^{h-1}(1-\eps_t)\leq\sum_t\eps_t$.
\end{lemma}
\begin{proof}
Start both processes in the same sampled state. As long as their states agree, use a maximal coupling of their next-state distributions, which preserves equality with probability at least $1-\eps_t$. The probability of no disagreement through time $h$ is at least the product of these factors. On the equality event terminal rewards are equal; on its complement their absolute difference is at most one. Taking expectations proves the first bound. The second is the union bound, or an induction on the product. Applying the argument separately to two branches bounds contrast distortion by the sum of their branch bounds. This bound is not generally sharp for a particular transition graph.
\end{proof}

\subsection{Endpoint detection boundary}
\begin{proof}[Proof of Theorem~\ref{thm:gap}]
First suppose $0<\gamma\leq2b$. Consider a null ideal pair $(\theta_R,\theta_F)=(1/2,1/2)$ and an alternative ideal pair $(1/2+\gamma/2,1/2-\gamma/2)$. Use the same observed means $(p_R,p_F)=(1/2+\gamma/4,1/2-\gamma/4)$ in both worlds, with independent Bernoulli arms and IID repetition. Each arm's distortion is $\gamma/4\leq b/2$. The entire observed data law is identical. If the test rejects with probability $u$ under this law, its type-I and type-II errors in these two worlds sum to $u+(1-u)=1$. They cannot both be at most $\delta<1/2$, irrespective of $n$. The construction uses valid Bernoulli parameters whenever $\gamma\leq1$; effects outside the feasible range are irrelevant.

Now suppose $g=\gamma-2b>0$. Let $W_i=Y_i^R-Y_i^F\in[-1,1]$ and $\widehat d=n^{-1}\sum_iW_i$. Under every null, $\E\widehat d\leq b$, and under every alternative $\E\widehat d\geq\gamma-b$. Threshold at $\gamma/2$. In either case the required one-sided deviation is at least $g/2$. Hoeffding's inequality for independent variables of range two gives each error at most $\exp(-ng^2/8)$. Thus the stated sufficient sample size follows. Dependence between arms within a pair does not affect this bound.

For the lower bound use independent arms. The null ideal pair is $(1/2,1/2)$ with observed pair
\[
 p^0=(1/2+b/2,\;1/2-b/2),
\]
and the alternative ideal pair is $(1/2+\gamma/2,1/2-\gamma/2)$ with observed pair
\[
 p^1=(1/2+(\gamma-b)/2,\;1/2-(\gamma-b)/2).
\]
All distortions have magnitude $b/2$. The difference between corresponding observed arm means has magnitude $g/2$. Under the stated parameter range the means lie in $[1/4,3/4]$. The Bernoulli inequality
$\kl(p,q)\leq(p-q)^2/[q(1-q)]$
therefore bounds each arm KL by $4g^2/3$. The two-arm KL is at most $8g^2/3\leq4g^2$, and $n$ IID pairs have KL at most $4ng^2$.

Let $E$ be the event that a proposed test rejects the null. Uniform error control implies $P_1(E)\geq1-\delta$ and $P_0(E)\leq\delta$. Data processing for KL applied to the binary indicator of $E$, followed by monotonicity of binary KL in this region, gives
\[
 4ng^2\geq\KL(P_1^{\otimes n},P_0^{\otimes n})
 \geq\kl(P_1(E),P_0(E))\geq\kl(1-\delta,\delta).
\]
Rearranging proves \eqref{eq:lower}. The logarithmic asymptotic dependence follows because $\kl(1-\delta,\delta)$ is of order $\log(1/\delta)$ for small $\delta$.
\end{proof}

\paragraph{Observed certification versus uniform power.}
If an observed population contrast is $d$, the intended effect lies in $[d-b,d+b]$ before endpoint clipping, so $d>b$ certifies positivity. In contrast, the worst-case observation of a true effect $\gamma$ is $\gamma-b$, while the largest observation compatible with a zero effect is $b$. Uniformly separating these sets requires $\gamma-b>b$. These are different quantifiers; replacing the testing boundary by $\gamma>b$ would be incorrect.

\subsection{Development-level confidence bounds}
\begin{proof}[Proof of Proposition~\ref{prop:certificate}]
For any one fixed contrast the variables $Z_i$ are independent, identically distributed, and in $[-1,1]$. Hoeffding's two-sided inequality yields
\[
 \Pr\{|\bar Z-\E Z_i|>r\}\leq2\exp(-nr^2/2).
\]
Set its right-hand side to $\delta/J$ and union bound over the $J$ contrasts to obtain \eqref{eq:hoeffding}. No independence between different contrasts is required.

For the second bound put $X_i=(Z_i+1)/2\in[0,1]$. Theorem~4 of \citet{maurer} gives a one-sided radius
$\sqrt{2s_X^2\log(2/\alpha)/n}+7\log(2/\alpha)/(3(n-1))$
with failure probability at most $\alpha$. Apply it to $X$ and $1-X$, each with $\alpha=\delta/(2J)$, and union bound over sides and contrasts. Since $s_X^2=s_Z^2/4$ and $Z=2X-1$, multiplication by two gives exactly \eqref{eq:eb}. The probability guarantee is fixed-sample; observing data until the bound happens to pass is not covered.

If the intended population contrast differs from $\E Z_i$ by at most a known $b$, expand either valid sampling interval by $b$ on each side. A per-development-pair distortion bound of at most $b$ implies this population bound by taking expectations. Intersecting with the known range $[-1,1]$ preserves coverage. If $b$ itself is obtained from a random calibration procedure, its failure probability needs a separate allocation and union bound; it is not free.
\end{proof}

\paragraph{Random exact identified intervals.}
There is an alternative when the kernels for each developed pair are exactly known and Theorem~\ref{thm:rectangular} supplies endpoints $\ell_i,u_i\in[-1,1]$. Since $\ell_i\leq\Delta_i\leq u_i$, averaging independent pairs and applying one-sided bounded-mean concentration to the two endpoint sequences yields a population interval. Estimating the kernels from data introduces another uncertainty layer. The paper's sensitivity experiments use known kernels and do not claim to have solved finite-data structural calibration.

\subsection{Nested variance, optimal allocation, and a lower bound}
\begin{proof}[Proof of Proposition~\ref{prop:variance}]
Conditional independence of the inner audits gives
$\Var(Z_i\mid\mathcal D_i)=\Var(W_{ij}\mid\mathcal D_i)/m$ and
$\E[Z_i\mid\mathcal D_i]=\Delta_i$.
The law of total variance yields $\Var(Z_i)=\tau^2+\sigma^2/m$. Averaging $n$ independent pairs divides this by $n$, proving \eqref{eq:variance}.

For continuous total budget $T=n(c_d+mc_a)$ substitute $n=T/(c_d+mc_a)$. Apart from $1/T$, the variance becomes
\[
 (\tau^2+\sigma^2/m)(c_d+mc_a)
 =\tau^2c_d+\sigma^2c_a+\tau^2c_am+\sigma^2c_d/m.
\]
Its derivative is $\tau^2c_a-\sigma^2c_d/m^2$, which vanishes at the stated square-root allocation. The second derivative is nonnegative, and the boundary $m\geq1$ gives the maximum with one. The continuous formula is feasible only when the total budget permits the corresponding $n$; a finite experiment compares neighboring integer allocations and includes a minimum number of outer replications. If $\tau^2=0$ or a cost vanishes, this interior formula does not apply.

For the replication lower bound, let $\Delta_i\in\{-a,+a\}$ for a fixed $a\in(0,1]$, with mixture probability $p$ for $+a$, and set every inner outcome deterministically to $W_{ij}=\Delta_i$. Arbitrarily many inner samples reveal only this one Bernoulli latent label per development pair. Take $p_0=1/2-u$, $p_1=1/2+u$, $0<u\leq1/4$, so the population effects differ by $4au$. An estimator accurate within $\eps$ with probability at least $3/4$ at both distributions would, for $u=\eps/a$ and $\eps\leq a/4$, produce a test by thresholding at zero; the means are $\pm2\eps$. Its two errors would be at most $1/4$. But the one-pair KL between the latent Bernoulli laws is at most $(2u)^2/[(1/2-u)(1/2+u)]\leq64u^2/3$. Data processing as above forces $n(64u^2/3)\geq\kl(3/4,1/4)>0$. Thus $n=\Omega(a^2/\eps^2)$ for fixed confidence, independent of the number of inner descendants.
\end{proof}

\subsection{Persistent adapter uncertainty is not rowwise uncertainty}\label{app:persistent}
The shared-history calculation in Corollary~\ref{cor:shared} concerns an uncertain initial distribution. A different kind of entanglement arises when an unknown \emph{persistent} adapter affects several transitions. This is important when the same interface translation is reused throughout a fork-and-transplant audit. We give a finite exact formulation; this is an application of the distinction between rectangular and nonrectangular uncertainty, not a new general robust-control algorithm.

Let $\mathcal K=\{1,\ldots,k\}$ be a finite collection of admissible adapters. Each $j\in\mathcal K$ specifies an entire sequence of transition matrices $K_{j,0},\ldots,K_{j,h-1}$ on a sufficient \emph{paired} state space. The state can include both branches and the common random-input state. It has known initial distribution $\nu$ and terminal contrast $d\in[-1,1]$. The adapter is selected once, before the audit, and remains fixed through all $h$ steps. Define
\[
 g_j=\nu^\top K_{j,0}K_{j,1}\cdots K_{j,h-1}d.
\]
\begin{proposition}[Persistent-adapter identification]\label{prop:persistent}
If the only uncertainty is an unknown fixed adapter, the exact identified set is $\{g_1,\ldots,g_k\}$. If arbitrary mixtures over adapters selected once per audit are allowed, the sharp identified interval is $[\min_j g_j,\max_j g_j]$. In contrast, independently choosing any row $K_{j,t}(s,\cdot)$ at each $(t,s)$ gives a rectangular outer interval. This outer interval can contain negative effects even when every admissible persistent adapter gives the same positive effect.
\end{proposition}
\begin{proof}
Conditional on fixed adapter $j$, the Markov product gives the expected terminal contrast $g_j$. With no mixture, these and only these values are possible. Under an arbitrary distribution $w$ on adapters, the contrast is $\sum_jw_jg_j$ by conditioning. The image of the probability simplex is the convex hull of the finitely many real numbers $g_j$, namely the stated interval; its endpoints and every intermediate value are attainable. Each fixed adapter induces a feasible collection in the rectangular row relaxation, so that relaxation contains the fixed-adapter effects and their interval hull. Equality does not generally hold because the relaxation may switch adapters within the same trajectory.

For strictness take two steps, two adapters $j\in\{0,1\}$, and states $s,m_0,m_1,c,e$. At the first step, adapter $j$ sends $s$ to $m_j$. At the second step it sends $m_j$ to $c$ and $m_{1-j}$ to $e$. Let $d(c)=a$ and $d(e)=-a$ for any $0<a\leq1$; other states have zero contrast. Every fixed adapter visits $c$, hence $g_0=g_1=a$. The persistent identified interval is $[a,a]$. A rowwise relaxation can use adapter 0 first and adapter 1 second, reaching $e$, or use consistent choices to reach $c$. Since no terminal contrast lies outside $[-a,a]$, the relaxed endpoints are exactly $[-a,a]$. These contrasts are realizable by rewards in $[0,1]$: use $r_R=(1+d)/2$ and $r_F=(1-d)/2$. Thus this is a paired audit, not an unbounded reward construction.
\end{proof}

The implemented five-state example uses $a=.2$ and returns $[.2,.2]$ versus $[-.2,.2]$. Independent tests enumerate rectangular row choices in additional small random models and compare them with the dynamic program. These checks establish the claimed distinction for the declared models; they do not estimate an unknown adapter uncertainty set in a real agent. An outer interval containing zero means the relaxation cannot certify the sign, not that an admissible persistent adapter actually has a negative effect.

\section{Experimental specification and supplementary results}\label{app:results}
\subsection{Exact finite-state development systems}
Let the incumbent mechanism state be $j\in\{0,1,2,3\}$, initially zero. A proposal $k$ is selected with probability $Q_{jk}$ under recursive authorship and $Q_{0k}$ under fixed authorship. The next incumbent is $\max(j,k)$. Define $q_0=(.75,.20,.04,.01)$ and
\[
 Q^+=\begin{pmatrix}
 .75&.20&.04&.01\\ .15&.35&.45&.05\\ .10&.10&.25&.55\\0&0&0&1
 \end{pmatrix},\quad
 Q^-=\begin{pmatrix}
 .75&.20&.04&.01\\ .95&.049&.001&0\\ .98&0&.019&.001\\0&0&0&1
 \end{pmatrix}.
\]
The zero-dividend system has all rows $q_0$. Its two development chains are identical under the paired random stream. For any of these systems, fixed-author transitions are $P^F_{jj}=\sum_{k\leq j}q_{0k}$ and $P^F_{jk}=q_{0k}$ for $k>j$; recursive transitions replace $q_0$ by $Q_j$. The downstream value $g_B(j)=1-(1-p_j)^B$ corresponds to $B$ independent success attempts, with $p=(.02,.04,.20,.30)$. Since this value increases in $j$, selection is monotone in downstream quality. Nevertheless the $Q^-$ rows become worse proposers, producing a negative recursion dividend after the first development step.

We evaluate exact values at $C\in\{1,2,4,8,16,32,64\}$ and $B\in\{1,4,16\}$. Each system has 4,096 independent random trajectories paired between its two regimes, and all 63 cells are retained. Trajectories are shared across checkpoints within a run, so cells are correlated; the finite-family bound does not require independent cells. All transition products, exact effects, Monte Carlo estimates, and decomposition terms are exported.

\subsection{Known-kernel sensitivity experiment}
States are clean-failure, clean-success, rare-failure, and rare-success. Initial law is $(.95,0,.05,0)$, reward is $(0,1,0,1)$, and horizon is eight. Each clean or rare failure transitions to success with hazard $.06$ for the recursive branch and $.02$ for the fixed branch. Nominal success is absorbing. Structural support permits transitions only within the same clean or rare block; under uncertainty, even a nominal success row may move to its block's failure state. This is intentional and included in the robust optimization.

In the rare-uncertainty condition only rows in the rare block have TV radius $\eps$. In the global condition every row does. Radii range from zero to $.20$ in increments of $.01$. The returned endpoint kernels attain the lower and upper contrasts to numerical precision. A separate invariant-scratch-state check gives terminal value $.7$ under arbitrarily large within-equivalence-class uncertainty. The common-history example of Corollary~\ref{cor:shared} is also exported, keeping that joint model distinct from independent-arm uncertainty.

\subsection{Statistical calibration experiments}
For the lineage experiment, draw $\Delta_i$ independently as $+.4$ or $-.4$ with equal probability. Conditional on its sign, an inner difference is that sign times an independent Bernoulli$(.4)$ variable. Thus $\tau^2=.16$, $\sigma^2=.24$, and the mean is zero. We use $n=16$ and $m\in\{1,4,16,64,256,1024\}$ over 10,000 repetitions. Naive pooled Hoeffding treats $nm$ differences as independent; the correct Hoeffding bound uses $n$ development averages. The Student-$t$ interval is descriptive and approximate, not finite-sample-valid for this discrete mixture. Its coverage also falls below 95\% at large $m$, illustrating why we do not substitute it for the stated confirmation certificate.

\begin{figure}[h]
\centering\includegraphics[width=.89\linewidth]{figures/lineage_coverage.pdf}
\caption{Coverage depends on the replication unit. The pooled-descendant interval becomes sharply anticonservative as inner replication increases; the distribution-free development-level interval is conservative. The approximate $t$ interval is not an exact certificate either.}
\end{figure}
\begin{figure}[h]
\centering\includegraphics[width=.89\linewidth]{figures/lineage_variance.pdf}
\caption{More descendants reduce the within-development term but leave the between-development floor. The exact variance is $(.16+.24/m)/16$.}
\end{figure}

For endpoint testing, set $b=.03$, $\gamma\in\{.04,.08,.12\}$, and $n\in\{64,256,1024,4096,16384\}$. When $\gamma\leq2b$, the two hypotheses use the identical midpoint observation law from Theorem~\ref{thm:gap}. Otherwise they use the least-separated laws in that proof. Each sample consists of independent Bernoulli arms, and the test thresholds the estimated contrast at $\gamma/2$. We retain both type-I and type-II rates over 10,000 repetitions. Below the boundary their sum stays near one; increasing $n$ cannot separate identical laws. Above it the exponential upper bound is conservative, as expected.

\subsection{Complete rule-table interpreter specification}
A circuit with $g$ gates has $3g+1$ integer genes. Each gate stores its operation and two source indices. Source indices must reference an input or an earlier gate. The final gene selects the output wire. The target always initially selects the last gate; mutations may change that output. All six-input truth tables are packed into unsigned 64-bit words and evaluated exactly by bitwise operations. An independent scalar interpreter checks random circuits on every input in the tests.

The policy is $m\in\{0,\ldots,4\}^6\times\{0,\ldots,3\}\times\{0,1\}$. The seed is $(2,2,2,2,2,2,0,0)$. A protected interpreter assigns circuit-gene tags 0 (operation), 1 (wire reference), or 2 (output). The probability of selecting gene $j$ is proportional to
\[
 2^{m_{\operatorname{tag}(j)}}2^{m_6\ind\{j\text{ is active}\}}.
\]
A gene is active when it belongs to the current output's backward dependency cone. The edit is a random legal reset, increment, or decrement modulo its cardinality, with weights $2^{m_3},2^{m_4},2^{m_5}$. A random reset is allowed to leave a gene unchanged; this still consumes a proposal and evaluation. The circuit accepts strictly better functional scores and accepts equal scores only when $m_7=1$. Its released score never decreases.

Policy genomes use tags $(0,0,0,1,1,1,2,2)$, and all policy genes are marked active. Recursive and fixed-author development apply the same generic mutation routine to this genome. Recursive authorship uses the current selected $m$; fixed authorship uses the seed as author while still editing the current selected $m$. Meta-level selection always requires a strictly better summed training score. Because all policy genes are active, the author's active multiplier cancels during policy mutation; its neutral-acceptance coordinate is used for circuit repair but not for meta-level selection. This limited overlap between task-improvement and meta-proposal parameters is part of the benchmark, not a general model of self-improvement. It may help explain small dividends and limits extrapolation.

The entire policy space has $5^6\cdot4\cdot2=125,000$ elements. Random search samples this whole space uniformly with replacement. It has the same number of policy evaluations as local recursive and fixed-author search, but can make larger policy jumps. This baseline is deliberately not removed when it performs better. The paper does not compare against every possible nonrecursive optimizer.

The task-family operation probabilities, ordered AND/OR/XOR/NAND, are $(.25,.25,.25,.25)$, $(.10,.10,.70,.10)$, and $(.35,.35,.05,.25)$. Targets with fewer than four or more than sixty true outputs are rejected. Five corruption attempts select genes from the target's initial active cone with replacement and randomly reset them. The active list is not recomputed during corruption; no-op corruptions are possible. A parent computing exactly the target function is rejected. The rejection sampler has an explicit maximum of 10,000 attempts per requested task. The reported population is this rejection-conditioned distribution, not uniform Boolean functions.

\subsection{Protocols, uncertainty families, and outcomes}
The initial protocol uses seeds 1000 through 1511 separately for each task family. Development streams use a seed sequence built from (seed, family index, 101); audit streams use (seed, family index, 202). Pilot seeds 0 through 3 are excluded. Six training tasks and their search random streams remain fixed within each development run. New audit task--parent pairs and search randomness are generated separately, then shared between paired methods. Candidate policy acceptance never uses final audit outcomes.

The independent confirmation uses seeds 100000 through 104095 per family, one checkpoint $C=128$, and one downstream budget $B=32$. The one-percentage-point margin was set after examining the initial study and before executing confirmation seeds. No policy, task-generator, or selection parameter was tuned between the studies. The original full sweep is retained rather than merged selectively with confirmation results.

\begin{figure}[h]
\centering\includegraphics[width=.89\linewidth]{figures/confirmation.pdf}
\caption{Independent confirmation with 4,096 development seeds per family. Simultaneous finite-sample intervals are entirely inside the prespecified $\pm1$ percentage-point band. This is a tolerance-specific conclusion, not proof of exact equality.}
\end{figure}

\begin{table}[h]
\centering\small
\caption{All nine seed-relative mechanism gains from independent confirmation. Empirical Bernstein intervals are simultaneous within this separate nine-comparison family at level 95\%. This is not a claim that the primary and secondary confidence families jointly have 95\% coverage.}
\begin{tabular}{llrr}
\toprule
Family & Method versus seed & Gain (pp) & Simultaneous 95\% interval (pp) \\
\midrule
Mixed & Recursive & 5.55 & $[4.58,6.51]$ \\
Mixed & Fixed author & 5.61 & $[4.65,6.58]$ \\
Mixed & Random search & 6.99 & $[6.06,7.91]$ \\
Xor & Recursive & 4.08 & $[3.12,5.04]$ \\
Xor & Fixed author & 4.10 & $[3.14,5.06]$ \\
Xor & Random search & 5.93 & $[4.99,6.88]$ \\
Logic & Recursive & 6.21 & $[5.24,7.19]$ \\
Logic & Fixed author & 6.15 & $[5.18,7.13]$ \\
Logic & Random search & 7.77 & $[6.83,8.71]$ \\
\bottomrule
\end{tabular}

\end{table}

The aggregate CSV additionally includes recursive-minus-random comparisons. Their variance-sensitive bounds were computed descriptively after confirmation and are labeled post hoc; they are not substituted for the prespecified primary analysis. Primary recursion-dividend intervals are based on $J=3$; seed-relative gain intervals use $J=9$. The original sweep retains its more conservative Hoeffding bounds ($J=27$ for recursive-minus-fixed cells and $J=135$ for the remaining summaries) and approximate $t$ intervals. No interval is presented as valid after arbitrary stopping.

\subsection{Cost accounting and reproducibility commands}
The code is single-process Python with NumPy, SciPy, and Numba. Initial execution used Python 3.13.5, NumPy 2.3.5, SciPy 1.17.0, and Numba 0.65.1 on a Linux CPU environment. Five logical CPUs were available, but no parallel Numba kernels, process pool, or GPU were used. The initial full suite took about 59.35 seconds and the independent confirmation about 314.28 seconds in that environment. These are measured runtimes, not guarantees for another laptop, and JIT compilation/cache effects influence them.

With six development tasks, repair budget 24, and $C$ policy proposals, each development regime consumes $(C+1)\cdot6\cdot25$ circuit evaluations. Audit of one frozen policy at budget $B$ consumes $32(B+1)$ circuit evaluations. Multiple audit budgets along the same repair trajectory share the work through the largest budget. Baseline seed audits are computed once per audit case and reused when forming differences. Additional final research scoring is not hidden inside one method's development budget.

To reproduce, install the pinned requirements, run the tests, then execute:
\begin{verbatim}
python scripts/run_experiments.py --suite all
python scripts/run_experiments.py --suite programs \
  --protocol protocol_confirmation.json \
  --out results/confirmation
python scripts/analyze.py
python scripts/plot_results.py
cd paper
pdflatex -interaction=nonstopmode -halt-on-error main.tex
pdflatex -interaction=nonstopmode -halt-on-error main.tex
\end{verbatim}
A separate smoke mode uses pilot seeds and writes outside the paper-result directory. Protocol hashes and source hashes are recorded with each run. Raw archives contain integer audit scores for every development seed, method, checkpoint, task, and budget, together with selected policies, training scores, acceptances, and initial parent scores. Exported per-run CSV files permit recomputing confidence intervals without rerunning search.

\section{Claim audit and remaining validation}\label{app:claims}
The exact counterexample, performance-difference identity, robust recursion, testing bounds, confidence formulas, and variance statements are proved above under their stated assumptions. Their general ingredients are established mathematics; the paper does not claim to invent robust dynamic programming, concentration inequalities, self-improving code, or descendant-based improver evaluation. Numerical tests independently compare TV optimization against linear programming, evaluate attained robust endpoints, check the telescoping identity, check scalar and packed Boolean semantics, and test noninterference of input arrays.

Important limitations remain for a submission decision. The relevance and novelty of the integrated causal-audit framing relative to existing RSI evaluation and robust-statistics work require a critical external reading. Full human verification of the mathematical claims and code has not yet been completed. The executable benchmark is narrow, and a second nontrivial family of improver representations would strengthen external validity. These are explicit validation needs, not unperformed experiments described as completed results. There are no pending theorem statements silently promoted to theorems in this manuscript.

\section{Conditional-tree replication and independent confirmation}\label{app:trees}
\subsection{Grammar, interpreter, and intervention}
The tree has three feature/threshold internal nodes in breadth-first order and four leaves. Its genotype is
\[
 (f_0,t_0,f_1,t_1,f_2,t_2,a_0,a_1,a_2,a_3),
 \quad f_j\in\{0,\ldots,4\},\ t_j\in\{0,\ldots,8\},\ a_j\in\{0,\ldots,23\}.
\]
At node $j$, compare feature $f_j$ with $t_j/8$; equality selects the right branch. After exactly two comparisons the selected leaf specifies an edit. Let $a=12u+3v+w$, where $u\in\{0,1\}$ restricts selection to active genes when one, $v\in\{0,1,2,3\}$ is a gene-type selector (3 means any type), and $w\in\{0,1,2\}$ chooses uniform replacement, modular increment, or modular decrement. Select uniformly among eligible coordinates; when that set is empty use all coordinates. This fallback is declared in advance and shared across regimes.

In circuits, gene types 0, 1, and 2 denote gate operations, wires, and the output pointer. In trees they denote feature selectors, thresholds, and leaf actions. This mapping is a declared interface design choice, not a uniquely correct mapping. All tree syntax positions are eligible as active at the meta level; circuit activity is determined by its output dependency cone. The tree features are: elapsed fraction of the declared horizon; stagnation clipped at eight and divided by eight; current score in $[0,1]$; active-gene fraction; and the mean genotype value after coordinate-wise division by cardinality minus one. At the meta level, current score is mean training accuracy across tasks.

The seed is $(0,2,1,4,2,4,9,9,9,9)$, whose leaves all implement uniform-coordinate random replacement. Recursive development uses the current selected tree to edit its own syntax, while fixed-author development applies the seed's editing behavior to the evolving incumbent syntax. Thus the initial proposal law is identical, and both regimes can discover better trees. Selection accepts only strict mean training-score improvement; the task-repair inner loop also accepts only strict score improvement. The original rule-table experiment instead allowed an editable neutral-acceptance option. A third baseline samples all ten tree genes independently and uniformly from their legal values.

\subsection{Sampling, budgets, and protocol chronology}
The Boolean task distribution is unchanged from the original code: six inputs, sixteen gates, nontrivial targets, and parents with five attempted corruptions on the target's active genes, rejecting functionally correct parents. Each development seed uses six training tasks with inner budget 24; every candidate including duplicates is reevaluated using fixed training randomness. Final audits use 24 fresh task--parent pairs with no shared development history. All 64 truth-table inputs are visible during inner-loop selection. We therefore test transfer across tasks, not generalization from hidden truth-table inputs.

For family index $f$ and run index $i$, five independent NumPy streams are spawned from \texttt{SeedSequence([master, f, i])}. They generate training tasks, audit tasks, training search randomness, meta-level randomness, and audit search randomness. Methods share the same appropriate random-input arrays, but never mutable state. Within a seed set, outcomes are paired, not independent; seed sets are the sampling units.

The exploratory protocol uses master seed 682719, 2,048 runs per family, development checkpoints 16 and 64, and audit budgets 16 and 64. The grammar and primary cell $(64,64)$ were fixed before these outcomes. The confirmation uses new master seed 682720, 4,096 runs per family, and only $(64,64)$. Confirmation was selected after the exploratory findings, with no change to grammar, seed policy, task family, optimization, or selection. Both protocols were committed before their respective outcomes; neither is described as external preregistration. Confidence families are separate across studies, and initial outcomes are not pooled into confirmation.

The tree's time feature makes the declared horizon part of the algorithm. Accordingly every audit budget is a separate restart, with identical random numbers used only for coupling. The exploratory development checkpoint at 16 is a prefix of a 64-step development schedule, not a separately re-timed 16-step algorithm. The main confirmation fixes both horizons at 64. All local methods have identical charged evaluation counts; control-flow and random-generation instruction costs need not coincide. Uniform whole-program search changes more syntax per proposal but pays the same candidate evaluation cost; this stronger baseline is disclosed rather than hidden.

\subsection{Primary outcomes and inferential units}
Each stored primary score is an integer sum over $24\times64=1,536$ correct truth-table outputs. Dividing the within-seed recursive-minus-fixed difference by 1,536 gives $Z_i\in[-1,1]$. The primary empirical Bernstein interval is Proposition~\ref{prop:certificate} with $n=4,096$, $J=3$, and $\delta=.05$. Nine seed-relative gains use a separate $J=9$ family. These are finite-sample simultaneous bounds, not normal approximations or bounds treating all outputs as independent. An additional, more conservative sensitivity check places all twelve confirmation contrasts in a single $J=12$ family; all three negative dividends and all nine positive seed gains retain their signs. This is not substituted for the originally specified primary intervals.

\begin{table}[h]
\centering
\caption{Conditional-tree confirmation seed gains in percentage points. All nine intervals share a separate simultaneous 95\% guarantee.}
\begin{tabular}{llrr}
\toprule
Family & Method & Gain & Simultaneous 95\% interval \\
\midrule
mixed & Recursive & 3.28 & [2.26, 4.31] \\
mixed & Fixed author & 4.78 & [3.76, 5.80] \\
mixed & Random search & 6.29 & [5.26, 7.31] \\
xor & Recursive & 2.89 & [1.85, 3.94] \\
xor & Fixed author & 4.23 & [3.19, 5.27] \\
xor & Random search & 5.90 & [4.84, 6.96] \\
logic & Recursive & 3.75 & [2.70, 4.80] \\
logic & Fixed author & 5.41 & [4.38, 6.44] \\
logic & Random search & 6.98 & [5.94, 8.01] \\
\bottomrule
\end{tabular}

\end{table}

The exploratory primary dividends were $-1.779$, $-1.298$, and $-1.486$ percentage points (mixed, XOR-heavy, logic-heavy). Their respective simultaneous intervals were $[-3.352,-.207]$, $[-2.847,.251]$, and $[-3.069,.096]$. Only the mixed interval excluded zero. In the independent confirmation, all three primary dividend intervals exclude zero (Table~\ref{tab:trees}). The absolute mean released accuracies for recursive, fixed-author, random, and seed policies are, respectively: mixed $(.76738,.78237,.79742,.73457)$; XOR-heavy $(.69076,.70412,.72084,.66184)$; logic-heavy $(.80529,.82191,.83758,.76780)$.

\begin{figure}[h]
\centering\includegraphics[width=.88\linewidth]{figures/tree_dividend.pdf}
\caption{Conditional-tree recursion dividends. Exploratory and independently confirmed estimates are displayed separately; error bars use the finite-sample simultaneous bounds within each study.}
\end{figure}

Post-hoc confirmation traces show that inherited authors make approximately 36--37 condition-changing proposals and 22--23 leaf-changing proposals per 64-step run, compared with 32--33 and 24--25 for fixed authors. They accept fewer improvements. These counts describe trajectories after the authorship intervention. Without additional mediator interventions, they do not prove that condition-edit preference caused the negative dividend. Whole-tree random search frequently changes both kinds of syntax; its diagnostic counts need not sum to the proposal count.

\subsection{Implementation checks and reproducibility}
The implementation includes independent scalar decision-tree traversal, a separately written unpacked Boolean evaluator and whole task-repair reference, typed mutation and input-immutability checks, deterministic replay, first-generation equality, strict selection monotonicity, no-op accounting, and exact per-budget restart tests. An extracted Boolean core is checked against the unchanged original task generator on all three families. These checks validate the stated software behavior, not the soundness of every scientific inference.

The initial tree execution used a local equivalent copy of the Boolean functions; before confirmation these functions were isolated into \texttt{src/boolean\_core.py} without changing task semantics. The first eight exploratory seed sets per family replayed bitwise across all recorded arrays after this import refactor. Exact executed initial sources, hashes, and the replay report are preserved. Confirmation hashes refer directly to the final factored implementation. The initial and confirmation executions took approximately 156 and 260 seconds, respectively, in the recorded single-process environment, including JIT compilation. These measurements are not runtime promises.

Commands to reproduce the separate studies are:
\begin{verbatim}
python scripts/run_trees.py --protocol protocol_trees.json \
  --output results/trees
python scripts/run_trees.py \
  --protocol protocol_trees_confirmation.json \
  --output results/trees_confirmation
python scripts/analyze_trees.py \
  --results results/trees_confirmation
python scripts/verify_trees.py
\end{verbatim}
The runner refuses to overwrite existing raw results. A full research archive contains integer per-development outcomes, all selected policies, training scores, acceptance counts, diagnostics, metadata, and source hashes. Repository summaries include exact integer sufficient statistics so the primary confidence bounds can also be checked without rerunning search. Initial and confirmation arrays remain separate. This second representation strengthens implementation diversity but is not a second application domain.

\section{A controlled diagnostic of the task-to-meta type mapping}\label{app:adapterdiagnostic}
The negative conditional-tree dividend raises a possible explanation: an editing preference useful for circuit operations and wires might map poorly to feature selectors, thresholds, and action leaves during self-editing. We test one precisely defined part of this explanation rather than infer it from post-treatment edit counts.

For each permutation $\pi$ of $\{0,1,2\}$, replace only the meta-level gene tag $v$ by $\pi(v)$. Keep circuit tags, genotype cardinalities, tree semantics, features, seed, evaluator, task distribution, and both budgets fixed. This changes the complete recursive system; it is not merely a semantics-preserving renaming. It is, however, a selective intervention with a useful invariance.

\begin{proposition}[Fixed-author invariance under meta-tag permutation]\label{prop:adapterinvariance}
For the declared seed with action 9 at every leaf, fixed-author development is pointwise identical under all six meta-tag permutations, for every training instance and random-input stream. Its final frozen-policy audit is therefore also identical. The same invariance applies to the first recursive proposal, but need not hold after a changed tree is inherited.
\end{proposition}
\begin{proof}
Action 9 selects all coordinates without inspecting tags or activity, then performs uniform replacement. All internal decisions of the seed reach a leaf with this same action. Starting from the same incumbent, every fixed-author proposal thus selects the same coordinate and replacement using the same two random numbers for every permutation. Evaluation uses unchanged task semantics and identical training randomness, so scores and acceptance decisions coincide. Induction on development steps gives identical incumbents, scores, and final mechanisms. Identical frozen mechanisms yield identical audits under a common audit input stream. Before the first recursive change, its author is the same seed, proving the final statement. A later inherited leaf can select a particular tag, so the induction does not apply to arbitrary recursive authors.
\end{proof}

An additional protocol, fixed after the primary confirmation but before these diagnostic outcomes, evaluates every permutation using 2,048 new independent seed sets per family (master seed 682721), $C=B=64$, six training tasks, and 24 audit tasks. Within each seed, six recursive variants are paired with one common fixed-author trajectory. Reusing that invariant baseline reduces redundant work, not the declared cost of any compared method. Each recursive-minus-identity contrast is a paired development-level observation in $[-1,1]$; the 15 nonidentity effects form one 95\% simultaneous empirical Bernstein family. The 18 recursive-minus-fixed dividends form a separate confidence family. Identity effects are exactly zero by construction. No permutation is selected for selective reporting.

\begin{table}[h]
\centering
\caption{All mapping diagnostics, in percentage points. A mapping 120 sends original tag 0 to 1, 1 to 2, and 2 to 0. Dividend and adapter-change intervals belong to separate simultaneous confidence families.}
{\small\begin{tabular}{llrr}
\toprule
Family & Mapping & $\Gamma$ [95\% interval] & Change from identity [95\% interval] \\
\midrule
mixed & 012 & -1.47 [-3.48, 0.54] & 0 (identity) \\
mixed & 021 & -1.35 [-3.36, 0.65] & 0.11 [-1.73, 1.96] \\
mixed & 102 & -1.46 [-3.47, 0.56] & 0.01 [-1.75, 1.77] \\
mixed & 120 & -1.35 [-3.36, 0.66] & 0.11 [-1.75, 1.98] \\
mixed & 201 & -1.34 [-3.35, 0.66] & 0.12 [-1.74, 1.99] \\
mixed & 210 & -1.34 [-3.34, 0.67] & 0.13 [-1.73, 1.99] \\
xor & 012 & -1.44 [-3.45, 0.57] & 0 (identity) \\
xor & 021 & -1.46 [-3.47, 0.56] & -0.01 [-1.84, 1.81] \\
xor & 102 & -1.44 [-3.46, 0.57] & -0.00 [-1.73, 1.73] \\
xor & 120 & -1.37 [-3.39, 0.64] & 0.07 [-1.79, 1.92] \\
xor & 201 & -1.49 [-3.50, 0.53] & -0.05 [-1.89, 1.79] \\
xor & 210 & -1.36 [-3.38, 0.65] & 0.08 [-1.77, 1.92] \\
logic & 012 & -1.63 [-3.67, 0.41] & 0 (identity) \\
logic & 021 & -1.61 [-3.65, 0.43] & 0.03 [-1.84, 1.89] \\
logic & 102 & -1.68 [-3.72, 0.37] & -0.04 [-1.83, 1.74] \\
logic & 120 & -1.52 [-3.55, 0.51] & 0.11 [-1.77, 1.99] \\
logic & 201 & -1.58 [-3.63, 0.46] & 0.05 [-1.84, 1.94] \\
logic & 210 & -1.49 [-3.52, 0.54] & 0.15 [-1.72, 2.01] \\
\bottomrule
\end{tabular}
}
\end{table}

All 18 point-estimated dividends are negative. Relative to identity, point-estimated changes range from approximately $-.05$ to $+.15$ percentage points, and every nonidentity adapter-effect interval includes zero. This diagnostic does not support a detected benefit from the tested remappings. However, its conservative intervals are too wide to establish practical equivalence among mappings, and the simultaneous dividend intervals in this smaller, larger-family study do not exclude zero. We therefore do not claim that all mappings have a certified negative dividend, or that this experiment disproves every task-to-meta mismatch explanation. The original independent confirmation retains its separate, prespecified inferential claim. The diagnostic took approximately 253 seconds on the recorded single-process CPU environment; selected policies and all raw outcomes are retained.

\end{document}
